\subsection{Restrições na Base}

O movimento dos elos do manipulador em ambiente de microgravidade gera reações na base, que correspondem ao corpo principal do \emph{chaser}.  
Essas reações - torques - devem ser limitados para que a atitude da base permaneça estável durante a operação de \emph{rendezvous and docking/berthing}, evitando que rotações excessivas ou desvios possam comprometer a missão.

As restrições são aplicadas à taxa angular da base e ao seu erro de orientação em relação a uma referência inercial ou ao alvo:

\begin{equation}
\label{eq:restr_base_w}
\|\vec{\omega}_B\| \leq \omega_{b,\max},
\end{equation}

\begin{equation}
\label{eq:restr_base_theta}
\angle(\mathbf{R}_B, \mathbf{R}_{B,\text{ref}}) \leq \Theta_{b,\max},
\end{equation}

onde:
\begin{itemize}
    \item $\vec{\omega}_B$: vetor de velocidade angular da base;
    \item $\omega_{b,\max}$: valor máximo admissível de velocidade angular;
    \item $\mathbf{R}_B$: matriz de rotação que descreve a atitude da base;
    \item $\mathbf{R}_{B,\text{ref}}$: matriz de rotação de referência (padrão de alinhamento desejado);
    \item $\Theta_{b,\max}$: limite máximo admissível para o desvio angular da base.
\end{itemize}

De forma prática, essas restrições asseguram que a base não sofra rotações acima de uma taxa controlável pelo sistema de controle de atitude, além disso, evita o desvio de orientação da base em relação à referência permaneça dentro de uma mergem segura, e que o manipulador opere sem induzir movimentos que comprometam a estabilidade global do manipulador robótico.

Os valores de $\omega_{b,\max}$ e $\Theta_{b,\max}$ serão definidos posteriormente com base em simulações de dinâmica acoplada entre base e manipulador, bem como nos requisitos de controle de atitude estabelecidos para a missão.
